% Generated draft: requires mathematical and source review.
\documentclass[11pt]{article}
\usepackage[margin=25mm]{geometry}
\usepackage{fontspec}
\setmainfont{Latin Modern Roman}
\usepackage{amsmath,amssymb,mathtools,unicode-math}
\setmathfont{Latin Modern Math}
\usepackage{enumitem,needspace,xurl,hyperref}
\setlength{\parindent}{0pt}
\setlength{\parskip}{5pt}
\setlength{\emergencystretch}{3em}
\newcommand{\N}{\mathbb{N}}
\newcommand{\Z}{\mathbb{Z}}
\newcommand{\Q}{\mathbb{Q}}
\newcommand{\R}{\mathbb{R}}
\newcommand{\C}{\mathbb{C}}
\newcommand{\F}{\mathbb{F}}
\newcommand{\A}{\mathbb{A}}
\newcommand{\PP}{\mathbb{P}}
\newcommand{\E}{\mathbb{E}}
\newcommand{\eps}{\varepsilon}
\newcommand{\dd}{\,\mathrm{d}}
\begin{document}
\section*{UnsolvedMath ID: 700711}
\textbf{Source problem code:} catalogue entry
\subsection{Definitions and mathematical statement}
For elements \(u,v\) of a group, use the commutator convention \([u,v]=u^{-1}v^{-1}uv\), and define the iterated commutators recursively by
\[
[u,{}_1v]=[u,v],\qquad [u,{}_{n+1}v]=[[u,{}_nv],v]\quad(n\geq 1).
\]
For positive integers \(a,b\), let
\[
G(a,b)=\left\langle x,y\ \middle|\ x=[x,{}_a y],\quad y=[y,{}_b x]\right\rangle.
\]
The question is whether
\[
\forall a,b\in\Z_{\geq 1},\qquad |G(a,b)|<\infty.
\]
The supplied record does not explicitly specify the parameter range or the commutator convention. The positive-integer range and the recursion above make the intended standard interpretation explicit, but require confirmation against the original source.

The source reports that \(G(1,b)=1\) and \(G(2,2)=1\), where \(1\) denotes the trivial group, and states that nothing is known about \(G(2,3)\). These are supplied source assertions, including a historical knowledge claim, and are not independently verified here.

The source further remarks conditionally that, if every minimal simple group were a quotient of some \(G(a,b)\), this would yield a sequence of words in two variables characterizing soluble groups. Here a quotient means the image of a surjective group homomorphism. In the customary finite-group usage, a minimal simple group is a finite nonabelian simple group whose proper subgroups are all soluble. A group is soluble if its derived series, obtained by repeatedly taking the subgroup generated by all commutators, reaches the trivial group after finitely many steps. The record does not specify the scope of the proposed characterization or the sequence of words; this remark is retained as motivation, not as an additional established conclusion.
\subsection{Short English statement}
R. Brandl asks whether every group \(G(a,b)\) given by the two displayed iterated-commutator relations is finite. The draft uses positive integer parameters and an explicit commutator convention, both requiring source confirmation. The reported special cases and the motivation involving minimal simple groups are retained as supplied assertions.
\subsection{Sources}
\begin{itemize}
\item Local UnsolvedMath registry, permanent ID \(700711\), exact source problem code \texttt{catalogue entry}, titled \emph{Kourovka Notebook Problem 11.18}; supplied complete source record, upstream revision \texttt{d371ee917ac9647fbb06cd6fc529deec6f3dda82}. The registry labels the problem \emph{open}; this label is supplied metadata, not an independently verified assessment.
\item R. Brandl, proposer, 1990. \emph{The Kourovka Notebook (Unsolved Problems in Group Theory)}, Problem 11.18, originating in the 11th issue (1990); the supplied background identifies the 21st-issue edition (2026), \texttt{arXiv:1401.0300v48}. The supplied collection description names E. I. Khukhro and V. D. Mazurov as editors. The original source has not been independently consulted here.
\item R. Brandl and J. S. Wilson, \emph{J. Algebra} \textbf{116} (1988), 334--341, cited in the supplied statement for the motivation concerning words characterizing soluble groups. No article title or link is supplied.
\end{itemize}
\end{document}
